\documentclass[12pt]{article}
\usepackage{amsmath, amssymb}
\usepackage[margin=1in]{geometry}
\setlength{\parskip}{6pt}
\title{QUBO Formulation: Max-Cut}
\date{}
\begin{document}
\maketitle

\subsection*{Max-Cut}

\paragraph{Problem Statement.}
Let $G=(V,E)$ be an undirected graph with vertex set $V$ and edge set $E$. Each edge $(i,j) \in E$ is associated with a non-negative weight $w_{ij}$. The Max-Cut problem seeks to partition the vertices $V$ into two disjoint subsets $V_0$ and $V_1$ such that the total weight of edges connecting a vertex in $V_0$ to a vertex in $V_1$ is maximized. Formally, we wish to find a partition that maximizes $\sum_{(i,j) \in E} w_{ij} \cdot \mathbb{I}(i \in V_0, j \in V_1 \lor i \in V_1, j \in V_0)$, where $\mathbb{I}(\cdot)$ is the indicator function.

\paragraph{Decision Variables.}
For each vertex $i \in V$, we introduce a binary decision variable $x_i \in \{0, 1\}$. The value $x_i = 0$ indicates that vertex $i$ is assigned to subset $V_0$, and $x_i = 1$ indicates assignment to subset $V_1$. The vector $x = (x_i)_{i \in V}$ fully specifies a partition of the graph.

\paragraph{QUBO Derivation.}
We derive the Quadratic Unconstrained Binary Optimization (QUBO) formulation by encoding the objective function into a Hamiltonian $H(x)$ to be minimized.

\textbf{Step 1: Objective formulation.}
For binary variables, the condition $x_i \neq x_j$ is equivalent to the algebraic expression $x_i + x_j - 2x_ix_j$. Thus, the weight of the cut contributed by edge $(i,j)$ is $w_{ij}(x_i + x_j - 2x_ix_j)$. The total cut weight is $\sum_{(i,j) \in E} w_{ij}(x_i + x_j - 2x_ix_j)$. Since QUBO solvers minimize energy, we define the Hamiltonian as the negative of the objective:
\begin{align}
H(x) = -\sum_{(i,j) \in E} w_{ij}(x_i + x_j - 2x_ix_j).
\end{align}
Expanding the terms yields:
\begin{align}
H(x) = \sum_{(i,j) \in E} \left( -w_{ij}x_i - w_{ij}x_j + 2w_{ij}x_ix_j \right).
\end{align}

\textbf{Step 2: Constraints.}
The problem specification contains no constraints; the objective is a pure maximization over the binary hypercube $\{0,1\}^{|V|}$. Consequently, there are no constraint terms to convert.

\textbf{Step 3: Penalty terms.}
As there are no constraints, no penalty terms are required. The Hamiltonian consists solely of the negated objective.

\textbf{Step 4: Combined Hamiltonian.}
The QUBO Hamiltonian is given directly by the expanded objective:
\begin{align}
H(x) = \sum_{(i,j) \in E} \left( 2w_{ij}x_ix_j - w_{ij}x_i - w_{ij}x_j \right).
\end{align}

\textbf{Step 5: Q matrix entries and offset.}
We collect coefficients for each variable $x_i$ and pair $x_ix_j$ to read off the QUBO parameters. The Hamiltonian takes the standard form $H(x) = \sum_{i \in V} Q_{i,i}x_i + \sum_{i,j \in V, i < j} Q_{i,j}x_ix_j + c$. The general closed-form expressions for the parameters are:
\begin{align}
Q_{i,i} &= -\sum_{j : (i,j) \in E} w_{ij}, \quad \forall i \in V, \\
Q_{i,j} &= 2w_{ij}, \quad \forall (i,j) \in E, \ i < j, \\
Q_{i,j} &= 0, \quad \forall (i,j) \notin E, \ i \neq j, \\
c &= 0.
\end{align}

\paragraph{Penalty Weight Justification.}
The problem formulation imposes no constraints; therefore, the QUBO requires no penalty terms and no penalty weights. The Hamiltonian directly encodes the optimization objective over the feasible domain $\{0,1\}^{|V|}$.

\paragraph{Correctness Argument.}
Minimizing $H(x)$ over $x \in \{0,1\}^{|V|}$ recovers the optimal Max-Cut solution. For any assignment $x$, the term $x_i + x_j - 2x_ix_j$ evaluates to $1$ if $x_i \neq x_j$ and $0$ otherwise. Thus, the original objective sums the weights of all edges crossing the cut defined by $x$. Since $H(x)$ is the negative of this sum, the global minimum of $H(x)$ corresponds to the partition that maximizes the total crossing weight. The absence of constraints ensures that all binary assignments are feasible, and the energy landscape is strictly determined by the cut quality, guaranteeing that the ground state of $H$ yields the optimal cut.

\paragraph{Illustrative Example.}
Consider a concrete instance with $N=3$ vertices $V=\{0,1,2\}$ forming a triangle graph where every pair of vertices is connected. The edge set is $E=\{(0,1), (1,2), (0,2)\}$ with weights $w_{01}, w_{12}, w_{02}$.

Instantiating the general formulas, the Hamiltonian becomes:
\begin{align}
H(x) &= \left( -w_{01}x_0 - w_{01}x_1 + 2w_{01}x_0x_1 \right) \nonumber \\
&\quad + \left( -w_{12}x_1 - w_{12}x_2 + 2w_{12}x_1x_2 \right) \nonumber \\
&\quad + \left( -w_{02}x_0 - w_{02}x_2 + 2w_{02}x_0x_2 \right).
\end{align}
Collecting terms for each variable and pair, the QUBO parameters are:
\begin{align}
Q_{0,0} &= -w_{01} - w_{02}, \\
Q_{1,1} &= -w_{01} - w_{12}, \\
Q_{2,2} &= -w_{12} - w_{02}, \\
Q_{0,1} &= 2w_{01}, \quad Q_{1,2} = 2w_{12}, \quad Q_{0,2} = 2w_{02}, \\
c &= 0.
\end{align}
The resulting Hamiltonian is:
\begin{align}
H(x) &= (-w_{01}-w_{02})x_0 + (-w_{01}-w_{12})x_1 + (-w_{12}-w_{02})x_2 \nonumber \\
&\quad + 2w_{01}x_0x_1 + 2w_{12}x_1x_2 + 2w_{02}x_0x_2.
\end{align}
For unit weights $w_{ij}=1$, the optimal cut includes two edges (energy $-2$). Configurations such as $x=(0,0,1)$, $x=(0,1,0)$, $x=(1,0,0)$, and their complements achieve this minimum energy, corresponding to partitions that separate one vertex from the other two.

\end{document}